\subsection{Corrigés des exercices}

\targetCorr{fsec1_ex1}
\begin{correctionbox}[Corrigé de l'exercice 1 : Profit d'un Long Call]
1. Calcul du Payoff à l'échéance : $\max(112 - 100, 0) = 12$ €.

2. Calcul du profit : Le profit correspond au Payoff moins la prime payée.
Profit $= 12 - 5 = 7$ €. L'opération est gagnante.
\navToExo{fsec1_ex1}
\end{correctionbox}

\targetCorr{fsec1_ex2}
\begin{correctionbox}[Corrigé de l'exercice 2 : Profit d'un Short Put]
1. À l'échéance, $S_T = 185$ et $K = 200$. L'acheteur du Put a intérêt à vendre à 200. Il exerce donc l'option.

2. Le vendeur du Put est forcé d'acheter à 200 € ce qui ne vaut que 185 €. Il subit une perte à l'exercice de $185 - 200 = -15$ €.

3. Son profit net intègre la prime reçue initialement (+10 €) : Profit $= 10 - 15 = -5$ €.
\navToExo{fsec1_ex2}
\end{correctionbox}

\targetCorr{fsec1_ex3}
\begin{correctionbox}[Corrigé de l'exercice 3 : Statut d'un Call (Moneyness)]
Pour un Call, l'option est ITM si $S > K$, ATM si $S=K$, OTM si $S < K$.
Le sous-jacent $S = 175$ \$.
\begin{itemize}[label=$\cdot$]
    \item $K_1 = 160$ \$ : $175 > 160 \implies$ Dans la monnaie (ITM).
    \item $K_2 = 175$ \$ : $175 = 175 \implies$ À la monnaie (ATM).
    \item $K_3 = 190$ \$ : $175 < 190 \implies$ Hors de la monnaie (OTM).
\end{itemize}
\navToExo{fsec1_ex3}
\end{correctionbox}

\targetCorr{fsec1_ex4}
\begin{correctionbox}[Corrigé de l'exercice 4 : Combinaison de flux]
Valeur finale du portefeuille = Valeur de l'action + Payoff du Put.
$V_T = S_T + \max(100 - S_T, 0)$.
\begin{itemize}[label=$\cdot$]
    \item Si $S_T \geq 100$ : Le Put vaut $0$. $V_T = S_T$.
    \item Si $S_T < 100$ : Le Put vaut $100 - S_T$. $V_T = S_T + (100 - S_T) = 100$.
\end{itemize}
En résumé, $V_T = \max(S_T, 100)$. Ce profil garantit un plancher à 100 tout en gardant tout le potentiel de hausse. Ce profil est équivalent à celui de l'achat d'un Call (à la constante du Strike près, justifiant la parité entre le Call et le Put).
\navToExo{fsec1_ex4}
\end{correctionbox}

\targetCorr{fsec1_ex5}
\begin{correctionbox}[Corrigé de l'exercice 5 : Stratégie Straddle]
Le Payoff global est la somme du Payoff du Call et du Payoff du Put.
\begin{itemize}[label=$\cdot$]
    \item Si $S_T = 120$ : Call = $\max(120 - 100, 0) = 20$. Put = $\max(100 - 120, 0) = 0$. Payoff total = $20$.
    \item Si $S_T = 80$ : Call = $0$. Put = $\max(100 - 80, 0) = 20$. Payoff total = $20$.
\end{itemize}
Dans les deux cas de forte volatilité, l'investisseur gagne 20, pour un coût initial de 10 (Profit net de +10).
\navToExo{fsec1_ex5}
\end{correctionbox}